% Options for packages loaded elsewhere
\PassOptionsToPackage{unicode}{hyperref}
\PassOptionsToPackage{hyphens}{url}
\documentclass[
]{article}
\usepackage{xcolor}
\usepackage[margin=2.54cm]{geometry}
\usepackage{amsmath,amssymb}
\setcounter{secnumdepth}{-\maxdimen} % remove section numbering
\usepackage{iftex}
\ifPDFTeX
  \usepackage[T1]{fontenc}
  \usepackage[utf8]{inputenc}
  \usepackage{textcomp} % provide euro and other symbols
\else % if luatex or xetex
  \usepackage{unicode-math} % this also loads fontspec
  \defaultfontfeatures{Scale=MatchLowercase}
  \defaultfontfeatures[\rmfamily]{Ligatures=TeX,Scale=1}
\fi
\usepackage{lmodern}
\ifPDFTeX\else
  % xetex/luatex font selection
\fi
% Use upquote if available, for straight quotes in verbatim environments
\IfFileExists{upquote.sty}{\usepackage{upquote}}{}
\IfFileExists{microtype.sty}{% use microtype if available
  \usepackage[]{microtype}
  \UseMicrotypeSet[protrusion]{basicmath} % disable protrusion for tt fonts
}{}
\makeatletter
\@ifundefined{KOMAClassName}{% if non-KOMA class
  \IfFileExists{parskip.sty}{%
    \usepackage{parskip}
  }{% else
    \setlength{\parindent}{0pt}
    \setlength{\parskip}{6pt plus 2pt minus 1pt}}
}{% if KOMA class
  \KOMAoptions{parskip=half}}
\makeatother
\usepackage{color}
\usepackage{fancyvrb}
\newcommand{\VerbBar}{|}
\newcommand{\VERB}{\Verb[commandchars=\\\{\}]}
\DefineVerbatimEnvironment{Highlighting}{Verbatim}{commandchars=\\\{\}}
% Add ',fontsize=\small' for more characters per line
\usepackage{framed}
\definecolor{shadecolor}{RGB}{248,248,248}
\newenvironment{Shaded}{\begin{snugshade}}{\end{snugshade}}
\newcommand{\AlertTok}[1]{\textcolor[rgb]{0.94,0.16,0.16}{#1}}
\newcommand{\AnnotationTok}[1]{\textcolor[rgb]{0.56,0.35,0.01}{\textbf{\textit{#1}}}}
\newcommand{\AttributeTok}[1]{\textcolor[rgb]{0.13,0.29,0.53}{#1}}
\newcommand{\BaseNTok}[1]{\textcolor[rgb]{0.00,0.00,0.81}{#1}}
\newcommand{\BuiltInTok}[1]{#1}
\newcommand{\CharTok}[1]{\textcolor[rgb]{0.31,0.60,0.02}{#1}}
\newcommand{\CommentTok}[1]{\textcolor[rgb]{0.56,0.35,0.01}{\textit{#1}}}
\newcommand{\CommentVarTok}[1]{\textcolor[rgb]{0.56,0.35,0.01}{\textbf{\textit{#1}}}}
\newcommand{\ConstantTok}[1]{\textcolor[rgb]{0.56,0.35,0.01}{#1}}
\newcommand{\ControlFlowTok}[1]{\textcolor[rgb]{0.13,0.29,0.53}{\textbf{#1}}}
\newcommand{\DataTypeTok}[1]{\textcolor[rgb]{0.13,0.29,0.53}{#1}}
\newcommand{\DecValTok}[1]{\textcolor[rgb]{0.00,0.00,0.81}{#1}}
\newcommand{\DocumentationTok}[1]{\textcolor[rgb]{0.56,0.35,0.01}{\textbf{\textit{#1}}}}
\newcommand{\ErrorTok}[1]{\textcolor[rgb]{0.64,0.00,0.00}{\textbf{#1}}}
\newcommand{\ExtensionTok}[1]{#1}
\newcommand{\FloatTok}[1]{\textcolor[rgb]{0.00,0.00,0.81}{#1}}
\newcommand{\FunctionTok}[1]{\textcolor[rgb]{0.13,0.29,0.53}{\textbf{#1}}}
\newcommand{\ImportTok}[1]{#1}
\newcommand{\InformationTok}[1]{\textcolor[rgb]{0.56,0.35,0.01}{\textbf{\textit{#1}}}}
\newcommand{\KeywordTok}[1]{\textcolor[rgb]{0.13,0.29,0.53}{\textbf{#1}}}
\newcommand{\NormalTok}[1]{#1}
\newcommand{\OperatorTok}[1]{\textcolor[rgb]{0.81,0.36,0.00}{\textbf{#1}}}
\newcommand{\OtherTok}[1]{\textcolor[rgb]{0.56,0.35,0.01}{#1}}
\newcommand{\PreprocessorTok}[1]{\textcolor[rgb]{0.56,0.35,0.01}{\textit{#1}}}
\newcommand{\RegionMarkerTok}[1]{#1}
\newcommand{\SpecialCharTok}[1]{\textcolor[rgb]{0.81,0.36,0.00}{\textbf{#1}}}
\newcommand{\SpecialStringTok}[1]{\textcolor[rgb]{0.31,0.60,0.02}{#1}}
\newcommand{\StringTok}[1]{\textcolor[rgb]{0.31,0.60,0.02}{#1}}
\newcommand{\VariableTok}[1]{\textcolor[rgb]{0.00,0.00,0.00}{#1}}
\newcommand{\VerbatimStringTok}[1]{\textcolor[rgb]{0.31,0.60,0.02}{#1}}
\newcommand{\WarningTok}[1]{\textcolor[rgb]{0.56,0.35,0.01}{\textbf{\textit{#1}}}}
\usepackage{graphicx}
\makeatletter
\newsavebox\pandoc@box
\newcommand*\pandocbounded[1]{% scales image to fit in text height/width
  \sbox\pandoc@box{#1}%
  \Gscale@div\@tempa{\textheight}{\dimexpr\ht\pandoc@box+\dp\pandoc@box\relax}%
  \Gscale@div\@tempb{\linewidth}{\wd\pandoc@box}%
  \ifdim\@tempb\p@<\@tempa\p@\let\@tempa\@tempb\fi% select the smaller of both
  \ifdim\@tempa\p@<\p@\scalebox{\@tempa}{\usebox\pandoc@box}%
  \else\usebox{\pandoc@box}%
  \fi%
}
% Set default figure placement to htbp
\def\fps@figure{htbp}
\makeatother
\setlength{\emergencystretch}{3em} % prevent overfull lines
\providecommand{\tightlist}{%
  \setlength{\itemsep}{0pt}\setlength{\parskip}{0pt}}
\usepackage{bookmark}
\IfFileExists{xurl.sty}{\usepackage{xurl}}{} % add URL line breaks if available
\urlstyle{same}
\hypersetup{
  pdftitle={Assignment 4: Data Wrangling (Spring 2026)},
  pdfauthor={Lex Clay},
  hidelinks,
  pdfcreator={LaTeX via pandoc}}

\title{Assignment 4: Data Wrangling (Spring 2026)}
\author{Lex Clay}
\date{}

\begin{document}
\maketitle

\subsection{OVERVIEW}\label{overview}

This exercise accompanies the lessons in Environmental Data Analytics on
Data Wrangling.\\
Do not use any AI tools in completing this assignment.

\subsection{Directions}\label{directions}

\begin{enumerate}
\def\labelenumi{\arabic{enumi}.}
\tightlist
\item
  Rename this file
  \texttt{\textless{}FirstLast\textgreater{}\_A04\_DataWrangling.Rmd}
  (replacing \texttt{\textless{}FirstLast\textgreater{}} with your first
  and last name).
\item
  Change ``Student Name'' on line 3 (above) with your name.
\item
  Work through the steps, \textbf{creating code and output} that fulfill
  each instruction.
\item
  Be sure to \textbf{answer the questions} in this assignment document.
\item
  When you have completed the assignment, \textbf{Knit} the text and
  code into a single PDF file.
\item
  \textbf{Ensure that code in code chunks is tidy and does not extend
  off the page in the PDF}.
\item
  \textbf{Push your completed RMD to your GitHub account}
\end{enumerate}

\subsection{Set up your session}\label{set-up-your-session}

1a. Load the \texttt{tidyverse} and \texttt{here} packages into your
session.

\begin{Shaded}
\begin{Highlighting}[]
\FunctionTok{library}\NormalTok{(tidyverse)}
\end{Highlighting}
\end{Shaded}

\begin{verbatim}
## -- Attaching core tidyverse packages ------------------------ tidyverse 2.0.0 --
## v dplyr     1.1.4     v readr     2.1.5
## v forcats   1.0.0     v stringr   1.5.2
## v ggplot2   4.0.0     v tibble    3.3.0
## v lubridate 1.9.4     v tidyr     1.3.1
## v purrr     1.1.0     
## -- Conflicts ------------------------------------------ tidyverse_conflicts() --
## x dplyr::filter() masks stats::filter()
## x dplyr::lag()    masks stats::lag()
## i Use the conflicted package (<http://conflicted.r-lib.org/>) to force all conflicts to become errors
\end{verbatim}

\begin{Shaded}
\begin{Highlighting}[]
\FunctionTok{library}\NormalTok{(here)}
\end{Highlighting}
\end{Shaded}

\begin{verbatim}
## Warning: package 'here' was built under R version 4.5.2
\end{verbatim}

\begin{verbatim}
## here() starts at C:/Users/lexcl/OneDrive/Documents/EDE/EDE_Spring2026
\end{verbatim}

\begin{Shaded}
\begin{Highlighting}[]
\FunctionTok{library}\NormalTok{(dplyr)}
\end{Highlighting}
\end{Shaded}

1b. Check your working directory.

\begin{Shaded}
\begin{Highlighting}[]
\FunctionTok{getwd}\NormalTok{()}
\end{Highlighting}
\end{Shaded}

\begin{verbatim}
## [1] "C:/Users/lexcl/OneDrive/Documents/EDE/EDE_Spring2026"
\end{verbatim}

\begin{Shaded}
\begin{Highlighting}[]
\FunctionTok{here}\NormalTok{()}
\end{Highlighting}
\end{Shaded}

\begin{verbatim}
## [1] "C:/Users/lexcl/OneDrive/Documents/EDE/EDE_Spring2026"
\end{verbatim}

1c. Read in all four raw data files associated with the EPA Air dataset,
being sure to set string columns to be read in a factors. See the README
file for the EPA air datasets for more information (especially if you
have not worked with air quality data previously).

\begin{enumerate}
\def\labelenumi{\arabic{enumi}.}
\setcounter{enumi}{1}
\tightlist
\item
  Add the appropriate code to reveal the dimensions of the four
  datasets.
\end{enumerate}

\begin{Shaded}
\begin{Highlighting}[]
\CommentTok{\#1a }

\CommentTok{\#1b }

\CommentTok{\#1c }
\NormalTok{EPAAirList }\OtherTok{\textless{}{-}} \FunctionTok{list}\NormalTok{(}
\NormalTok{  EPAAirOne }\OtherTok{\textless{}{-}} \FunctionTok{read.csv}\NormalTok{(}\AttributeTok{file=}\FunctionTok{here}\NormalTok{(}\StringTok{\textquotesingle{}./Data/Raw/EPAair\_O3\_NC2018\_raw.csv\textquotesingle{}}\NormalTok{),}\AttributeTok{stringsAsFactors=}\ConstantTok{TRUE}\NormalTok{),}
\NormalTok{  EPAAirTwo }\OtherTok{\textless{}{-}}\FunctionTok{read.csv}\NormalTok{(}\AttributeTok{file=}\FunctionTok{here}\NormalTok{(}\StringTok{\textquotesingle{}./Data/Raw/EPAair\_O3\_NC2019\_raw.csv\textquotesingle{}}\NormalTok{),}\AttributeTok{stringsAsFactors=}\ConstantTok{TRUE}\NormalTok{),}
\NormalTok{  EPAAirThree }\OtherTok{\textless{}{-}} \FunctionTok{read.csv}\NormalTok{(}\AttributeTok{file=}\FunctionTok{here}\NormalTok{(}\StringTok{\textquotesingle{}./Data/Raw/EPAair\_PM25\_NC2018\_raw.csv\textquotesingle{}}\NormalTok{),}\AttributeTok{stringsAsFactors=}\ConstantTok{TRUE}\NormalTok{),}
\NormalTok{  EPAAirFour }\OtherTok{\textless{}{-}} \FunctionTok{read.csv}\NormalTok{(}\AttributeTok{file=}\FunctionTok{here}\NormalTok{(}\StringTok{\textquotesingle{}./Data/Raw/EPAair\_PM25\_NC2019\_raw.csv\textquotesingle{}}\NormalTok{),}\AttributeTok{stringsAsFactors=}\ConstantTok{TRUE}\NormalTok{)}
\NormalTok{)}
    
\FunctionTok{str}\NormalTok{(EPAAirList)}
\end{Highlighting}
\end{Shaded}

\begin{verbatim}
## List of 4
##  $ :'data.frame':    9737 obs. of  20 variables:
##   ..$ Date                                : Factor w/ 364 levels "01/01/2018","01/02/2018",..: 60 61 62 63 64 65 66 67 68 69 ...
##   ..$ Source                              : Factor w/ 1 level "AQS": 1 1 1 1 1 1 1 1 1 1 ...
##   ..$ Site.ID                             : int [1:9737] 370030005 370030005 370030005 370030005 370030005 370030005 370030005 370030005 370030005 370030005 ...
##   ..$ POC                                 : int [1:9737] 1 1 1 1 1 1 1 1 1 1 ...
##   ..$ Daily.Max.8.hour.Ozone.Concentration: num [1:9737] 0.043 0.046 0.047 0.049 0.047 0.03 0.036 0.044 0.049 0.043 ...
##   ..$ UNITS                               : Factor w/ 1 level "ppm": 1 1 1 1 1 1 1 1 1 1 ...
##   ..$ DAILY_AQI_VALUE                     : int [1:9737] 40 43 44 45 44 28 33 41 45 40 ...
##   ..$ Site.Name                           : Factor w/ 40 levels "","Beaufort",..: 35 35 35 35 35 35 35 35 35 35 ...
##   ..$ DAILY_OBS_COUNT                     : int [1:9737] 17 17 17 17 17 17 17 17 17 17 ...
##   ..$ PERCENT_COMPLETE                    : num [1:9737] 100 100 100 100 100 100 100 100 100 100 ...
##   ..$ AQS_PARAMETER_CODE                  : int [1:9737] 44201 44201 44201 44201 44201 44201 44201 44201 44201 44201 ...
##   ..$ AQS_PARAMETER_DESC                  : Factor w/ 1 level "Ozone": 1 1 1 1 1 1 1 1 1 1 ...
##   ..$ CBSA_CODE                           : int [1:9737] 25860 25860 25860 25860 25860 25860 25860 25860 25860 25860 ...
##   ..$ CBSA_NAME                           : Factor w/ 17 levels "","Asheville, NC",..: 9 9 9 9 9 9 9 9 9 9 ...
##   ..$ STATE_CODE                          : int [1:9737] 37 37 37 37 37 37 37 37 37 37 ...
##   ..$ STATE                               : Factor w/ 1 level "North Carolina": 1 1 1 1 1 1 1 1 1 1 ...
##   ..$ COUNTY_CODE                         : int [1:9737] 3 3 3 3 3 3 3 3 3 3 ...
##   ..$ COUNTY                              : Factor w/ 32 levels "Alexander","Avery",..: 1 1 1 1 1 1 1 1 1 1 ...
##   ..$ SITE_LATITUDE                       : num [1:9737] 35.9 35.9 35.9 35.9 35.9 ...
##   ..$ SITE_LONGITUDE                      : num [1:9737] -81.2 -81.2 -81.2 -81.2 -81.2 ...
##  $ :'data.frame':    10592 obs. of  20 variables:
##   ..$ Date                                : Factor w/ 365 levels "01/01/2019","01/02/2019",..: 1 2 3 4 5 6 7 8 9 10 ...
##   ..$ Source                              : Factor w/ 2 levels "AirNow","AQS": 1 1 1 1 1 1 1 1 1 1 ...
##   ..$ Site.ID                             : int [1:10592] 370030005 370030005 370030005 370030005 370030005 370030005 370030005 370030005 370030005 370030005 ...
##   ..$ POC                                 : int [1:10592] 1 1 1 1 1 1 1 1 1 1 ...
##   ..$ Daily.Max.8.hour.Ozone.Concentration: num [1:10592] 0.029 0.018 0.016 0.022 0.037 0.037 0.029 0.038 0.038 0.03 ...
##   ..$ UNITS                               : Factor w/ 1 level "ppm": 1 1 1 1 1 1 1 1 1 1 ...
##   ..$ DAILY_AQI_VALUE                     : int [1:10592] 27 17 15 20 34 34 27 35 35 28 ...
##   ..$ Site.Name                           : Factor w/ 38 levels "","Beaufort",..: 33 33 33 33 33 33 33 33 33 33 ...
##   ..$ DAILY_OBS_COUNT                     : int [1:10592] 24 24 24 24 24 24 24 24 24 24 ...
##   ..$ PERCENT_COMPLETE                    : num [1:10592] 100 100 100 100 100 100 100 100 100 100 ...
##   ..$ AQS_PARAMETER_CODE                  : int [1:10592] 44201 44201 44201 44201 44201 44201 44201 44201 44201 44201 ...
##   ..$ AQS_PARAMETER_DESC                  : Factor w/ 1 level "Ozone": 1 1 1 1 1 1 1 1 1 1 ...
##   ..$ CBSA_CODE                           : int [1:10592] 25860 25860 25860 25860 25860 25860 25860 25860 25860 25860 ...
##   ..$ CBSA_NAME                           : Factor w/ 15 levels "","Asheville, NC",..: 8 8 8 8 8 8 8 8 8 8 ...
##   ..$ STATE_CODE                          : int [1:10592] 37 37 37 37 37 37 37 37 37 37 ...
##   ..$ STATE                               : Factor w/ 1 level "North Carolina": 1 1 1 1 1 1 1 1 1 1 ...
##   ..$ COUNTY_CODE                         : int [1:10592] 3 3 3 3 3 3 3 3 3 3 ...
##   ..$ COUNTY                              : Factor w/ 30 levels "Alexander","Avery",..: 1 1 1 1 1 1 1 1 1 1 ...
##   ..$ SITE_LATITUDE                       : num [1:10592] 35.9 35.9 35.9 35.9 35.9 ...
##   ..$ SITE_LONGITUDE                      : num [1:10592] -81.2 -81.2 -81.2 -81.2 -81.2 ...
##  $ :'data.frame':    8983 obs. of  20 variables:
##   ..$ Date                          : Factor w/ 365 levels "01/01/2018","01/02/2018",..: 2 5 8 11 14 17 20 23 26 29 ...
##   ..$ Source                        : Factor w/ 1 level "AQS": 1 1 1 1 1 1 1 1 1 1 ...
##   ..$ Site.ID                       : int [1:8983] 370110002 370110002 370110002 370110002 370110002 370110002 370110002 370110002 370110002 370110002 ...
##   ..$ POC                           : int [1:8983] 1 1 1 1 1 1 1 1 1 1 ...
##   ..$ Daily.Mean.PM2.5.Concentration: num [1:8983] 2.9 3.7 5.3 0.8 2.5 4.5 1.8 2.5 4.2 1.7 ...
##   ..$ UNITS                         : Factor w/ 1 level "ug/m3 LC": 1 1 1 1 1 1 1 1 1 1 ...
##   ..$ DAILY_AQI_VALUE               : int [1:8983] 12 15 22 3 10 19 8 10 18 7 ...
##   ..$ Site.Name                     : Factor w/ 25 levels "","Blackstone",..: 15 15 15 15 15 15 15 15 15 15 ...
##   ..$ DAILY_OBS_COUNT               : int [1:8983] 1 1 1 1 1 1 1 1 1 1 ...
##   ..$ PERCENT_COMPLETE              : num [1:8983] 100 100 100 100 100 100 100 100 100 100 ...
##   ..$ AQS_PARAMETER_CODE            : int [1:8983] 88502 88502 88502 88502 88502 88502 88502 88502 88502 88502 ...
##   ..$ AQS_PARAMETER_DESC            : Factor w/ 2 levels "Acceptable PM2.5 AQI & Speciation Mass",..: 1 1 1 1 1 1 1 1 1 1 ...
##   ..$ CBSA_CODE                     : int [1:8983] NA NA NA NA NA NA NA NA NA NA ...
##   ..$ CBSA_NAME                     : Factor w/ 14 levels "","Asheville, NC",..: 1 1 1 1 1 1 1 1 1 1 ...
##   ..$ STATE_CODE                    : int [1:8983] 37 37 37 37 37 37 37 37 37 37 ...
##   ..$ STATE                         : Factor w/ 1 level "North Carolina": 1 1 1 1 1 1 1 1 1 1 ...
##   ..$ COUNTY_CODE                   : int [1:8983] 11 11 11 11 11 11 11 11 11 11 ...
##   ..$ COUNTY                        : Factor w/ 21 levels "Avery","Buncombe",..: 1 1 1 1 1 1 1 1 1 1 ...
##   ..$ SITE_LATITUDE                 : num [1:8983] 36 36 36 36 36 ...
##   ..$ SITE_LONGITUDE                : num [1:8983] -81.9 -81.9 -81.9 -81.9 -81.9 ...
##  $ :'data.frame':    8581 obs. of  7 variables:
##   ..$ Date              : Factor w/ 365 levels "2019-01-01","2019-01-02",..: 3 6 9 12 15 18 21 24 27 30 ...
##   ..$ Daily_AQI_Value   : int [1:8581] 7 4 5 26 11 5 6 6 15 7 ...
##   ..$ Site_Name         : Factor w/ 25 levels "","Board Of Ed. Bldg.",..: 14 14 14 14 14 14 14 14 14 14 ...
##   ..$ AQS_Parameter_Desc: Factor w/ 2 levels "Acceptable PM2.5 AQI & Speciation Mass",..: 1 1 1 1 1 1 1 1 1 1 ...
##   ..$ County            : Factor w/ 21 levels "Avery","Buncombe",..: 1 1 1 1 1 1 1 1 1 1 ...
##   ..$ Site_Latitude     : num [1:8581] 36 36 36 36 36 ...
##   ..$ Site_Longitude    : num [1:8581] -81.9 -81.9 -81.9 -81.9 -81.9 ...
\end{verbatim}

\begin{Shaded}
\begin{Highlighting}[]
\CommentTok{\#dim(EPAAirList) this returns null when used with list. only works with single objects? }

\CommentTok{\#2 }
\end{Highlighting}
\end{Shaded}

\textbf{TIP}: All four datasets should have the same number of columns
but unique record counts (rows). Do your datasets follow this pattern?

\subsection{Wrangle individual datasets to create processed
files.}\label{wrangle-individual-datasets-to-create-processed-files.}

\begin{enumerate}
\def\labelenumi{\arabic{enumi}.}
\setcounter{enumi}{2}
\item
  Change any date columns to be date objects.
\item
  Create new dataframes with just the following columns:\\
  `Date', `DAILY\_AQI\_VALUE', `Site.Name', `AQS\_PARAMETER\_DESC',
  `COUNTY', `SITE\_LATITUDE', `SITE\_LONGITUDE'
\item
  For the PM2.5 datasets, fill all cells in AQS\_PARAMETER\_DESC with
  ``PM2.5'' (all cell values in this column should be identical).
\item
  Save all four processed datasets in the Processed folder. Use the same
  file names as the raw files but replace ``raw'' with ``processed''.
\end{enumerate}

\begin{Shaded}
\begin{Highlighting}[]
\CommentTok{\#3}
\NormalTok{EPAAirOne}\SpecialCharTok{$}\NormalTok{Date }\OtherTok{\textless{}{-}} \FunctionTok{mdy}\NormalTok{(EPAAirOne}\SpecialCharTok{$}\NormalTok{Date)}
\end{Highlighting}
\end{Shaded}

\begin{Shaded}
\begin{Highlighting}[]
\NormalTok{EPAAirTwo}\SpecialCharTok{$}\NormalTok{Date }\OtherTok{\textless{}{-}}\FunctionTok{mdy}\NormalTok{ (EPAAirTwo}\SpecialCharTok{$}\NormalTok{Date)}
\end{Highlighting}
\end{Shaded}

\begin{Shaded}
\begin{Highlighting}[]
\NormalTok{EPAAirThree}\SpecialCharTok{$}\NormalTok{Date }\OtherTok{\textless{}{-}} \FunctionTok{mdy}\NormalTok{(EPAAirThree}\SpecialCharTok{$}\NormalTok{Date)}
\end{Highlighting}
\end{Shaded}

\begin{Shaded}
\begin{Highlighting}[]
\NormalTok{EPAAirFour}\SpecialCharTok{$}\NormalTok{Date }\OtherTok{\textless{}{-}}\FunctionTok{ymd}\NormalTok{(EPAAirFour}\SpecialCharTok{$}\NormalTok{Date)}
\end{Highlighting}
\end{Shaded}

\#4

\begin{Shaded}
\begin{Highlighting}[]
\NormalTok{  EPAOne\_DataFrame }\OtherTok{\textless{}{-}} \FunctionTok{data.frame}\NormalTok{(}
    \AttributeTok{Date =}\NormalTok{ EPAAirOne}\SpecialCharTok{$}\NormalTok{Date,}
    \AttributeTok{Daily\_AQI\_Value=}\NormalTok{ EPAAirOne}\SpecialCharTok{$}\NormalTok{DAILY\_AQI\_VALUE, }
    \AttributeTok{Site\_Name =}\NormalTok{ EPAAirOne}\SpecialCharTok{$}\NormalTok{Site.Name, }
    \AttributeTok{AQS\_Parameter\_Desc =}\NormalTok{ EPAAirOne}\SpecialCharTok{$}\NormalTok{AQS\_PARAMETER\_DESC, }
    \AttributeTok{County =}\NormalTok{ EPAAirOne}\SpecialCharTok{$}\NormalTok{COUNTY, }
    \AttributeTok{Site\_Latitude =}\NormalTok{ EPAAirOne}\SpecialCharTok{$}\NormalTok{SITE\_LATITUDE, }
    \AttributeTok{Site\_Longitude =}\NormalTok{ EPAAirOne}\SpecialCharTok{$}\NormalTok{SITE\_LONGITUDE,}
    \AttributeTok{stringsAsFactors =} \ConstantTok{TRUE}
\NormalTok{  )}

\NormalTok{ EPATwo\_DataFrame }\OtherTok{\textless{}{-}} \FunctionTok{data.frame}\NormalTok{(}
    \AttributeTok{Date =}\NormalTok{ EPAAirTwo}\SpecialCharTok{$}\NormalTok{Date,}
    \AttributeTok{Daily\_AQI\_Value=}\NormalTok{ EPAAirTwo}\SpecialCharTok{$}\NormalTok{DAILY\_AQI\_VALUE, }
    \AttributeTok{Site\_Name =}\NormalTok{ EPAAirTwo}\SpecialCharTok{$}\NormalTok{Site.Name, }
    \AttributeTok{AQS\_Parameter\_Desc =}\NormalTok{ EPAAirTwo}\SpecialCharTok{$}\NormalTok{AQS\_PARAMETER\_DESC, }
    \AttributeTok{County =}\NormalTok{ EPAAirTwo}\SpecialCharTok{$}\NormalTok{COUNTY, }
    \AttributeTok{Site\_Latitude =}\NormalTok{ EPAAirTwo}\SpecialCharTok{$}\NormalTok{SITE\_LATITUDE, }
    \AttributeTok{Site\_Longitude =}\NormalTok{ EPAAirTwo}\SpecialCharTok{$}\NormalTok{SITE\_LONGITUDE, }
    \AttributeTok{stringsAsFactors =} \ConstantTok{TRUE}
   
\NormalTok{ )}
 
\NormalTok{ EPAThree\_DataFrame }\OtherTok{\textless{}{-}} \FunctionTok{data.frame}\NormalTok{(}
    \AttributeTok{Date =}\NormalTok{ EPAAirThree}\SpecialCharTok{$}\NormalTok{Date,}
    \AttributeTok{Daily\_AQI\_Value=}\NormalTok{ EPAAirThree}\SpecialCharTok{$}\NormalTok{DAILY\_AQI\_VALUE, }
    \AttributeTok{Site\_Name =}\NormalTok{ EPAAirThree}\SpecialCharTok{$}\NormalTok{Site.Name, }
    \AttributeTok{AQS\_Parameter\_Desc =}\NormalTok{ EPAAirThree}\SpecialCharTok{$}\NormalTok{AQS\_PARAMETER\_DESC, }
    \AttributeTok{County =}\NormalTok{ EPAAirThree}\SpecialCharTok{$}\NormalTok{COUNTY, }
    \AttributeTok{Site\_Latitude =}\NormalTok{ EPAAirThree}\SpecialCharTok{$}\NormalTok{SITE\_LATITUDE, }
    \AttributeTok{Site\_Longitude =}\NormalTok{ EPAAirThree}\SpecialCharTok{$}\NormalTok{SITE\_LONGITUDE, }
    \AttributeTok{stringsAsFactors =} \ConstantTok{TRUE}
\NormalTok{ )}
 
\NormalTok{ EPAFour\_DataFrame }\OtherTok{\textless{}{-}}\NormalTok{ EPAAirFour }\CommentTok{\#I think I somehow overrode the raw file with the processed one }
 \CommentTok{\#and I don\textquotesingle{}t know how to fix this!}
\end{Highlighting}
\end{Shaded}

\#5

\begin{Shaded}
\begin{Highlighting}[]
\NormalTok{EPAThreee\_DataFrame }\OtherTok{\textless{}{-}}\NormalTok{ EPAThree\_DataFrame }\SpecialCharTok{\%\textgreater{}\%} 
  \FunctionTok{mutate}\NormalTok{(}\AttributeTok{AQS\_Parameter\_Desc =} \StringTok{"PM2.5"}\NormalTok{)}



\NormalTok{EPAFourr\_DataFrame }\OtherTok{\textless{}{-}}\NormalTok{ EPAFour\_DataFrame}\SpecialCharTok{\%\textgreater{}\%} 
  \FunctionTok{mutate}\NormalTok{(}\AttributeTok{AQS\_Parameter\_Desc =} \StringTok{"PM2.5"}\NormalTok{)}
\end{Highlighting}
\end{Shaded}

\#6

\begin{Shaded}
\begin{Highlighting}[]
\CommentTok{\#Saving to processed csvs}

\FunctionTok{write.csv}\NormalTok{(}
\NormalTok{  EPAOne\_DataFrame, }
  \AttributeTok{row.names =} \ConstantTok{FALSE}\NormalTok{, }
  \AttributeTok{file=}\FunctionTok{here}\NormalTok{(}\StringTok{\textquotesingle{}./Data/Processed/EPAair\_03\_NC2018\_processed.csv\textquotesingle{}}\NormalTok{)}
\NormalTok{)}

\FunctionTok{write.csv}\NormalTok{(}
\NormalTok{  EPATwo\_DataFrame, }
  \AttributeTok{row.names =} \ConstantTok{FALSE}\NormalTok{, }
  \AttributeTok{file=}\FunctionTok{here}\NormalTok{(}\StringTok{\textquotesingle{}./Data/processed/EPAair\_O3\_NC2019\_processed.csv\textquotesingle{}}\NormalTok{)}
\NormalTok{)}

\FunctionTok{write.csv}\NormalTok{(}
\NormalTok{  EPAThreee\_DataFrame,}
  \AttributeTok{row.names=} \ConstantTok{FALSE}\NormalTok{, }
  \AttributeTok{file=}\FunctionTok{here}\NormalTok{(}\StringTok{\textquotesingle{}./Data/Processed/EPAair\_PM25\_NC2018\_processed.csv\textquotesingle{}}\NormalTok{)}
\NormalTok{)}

\FunctionTok{write.csv}\NormalTok{(}
\NormalTok{  EPAFourr\_DataFrame,}
  \AttributeTok{row.names =} \ConstantTok{FALSE}\NormalTok{, }
  \AttributeTok{file=}\FunctionTok{here}\NormalTok{(}\StringTok{\textquotesingle{}./Data/Processed/EPAair\_PM25\_NC2019\_processed.csv\textquotesingle{}}\NormalTok{)}
\NormalTok{)}
\end{Highlighting}
\end{Shaded}

\subsection{Combine datasets}\label{combine-datasets}

\begin{enumerate}
\def\labelenumi{\arabic{enumi}.}
\setcounter{enumi}{6}
\tightlist
\item
  Combine the four datasets with \texttt{rbind}. Make sure your column
  names are identical prior to running this code.'
\end{enumerate}

\begin{Shaded}
\begin{Highlighting}[]
\NormalTok{Master\_EPA\_Air }\OtherTok{\textless{}{-}} \FunctionTok{rbind}\NormalTok{(EPAOne\_DataFrame, EPATwo\_DataFrame, EPAThreee\_DataFrame, EPAFourr\_DataFrame)}
\end{Highlighting}
\end{Shaded}

\begin{enumerate}
\def\labelenumi{\arabic{enumi}.}
\setcounter{enumi}{7}
\tightlist
\item
  Use code to display the dimensions of the combined dataset.
\end{enumerate}

\begin{Shaded}
\begin{Highlighting}[]
\FunctionTok{dim}\NormalTok{(Master\_EPA\_Air)}
\end{Highlighting}
\end{Shaded}

\begin{verbatim}
## [1] 37893     7
\end{verbatim}

\begin{enumerate}
\def\labelenumi{\arabic{enumi}.}
\setcounter{enumi}{8}
\tightlist
\item
  Wrangle your new dataset with a pipe function (\%\textgreater\%) so
  that it fills the following conditions:
\end{enumerate}

\begin{itemize}
\item
  Include only sites that the four data frames have in common:\\
  ``Linville Falls'', ``Durham Armory'', ``Leggett'', ``Hattie
  Avenue'',\\
  ``Clemmons Middle'', ``Mendenhall School'', ``Frying Pan Mountain'',
  ``West Johnston Co.'', ``Garinger High School'', ``Castle Hayne'',
  ``Pitt Agri. Center'', ``Bryson City'', ``Millbrook School''
\item
  Some sites have multiple measurements per day. Use the
  split-apply-combine strategy to generate daily means: group by date,
  site name, AQS parameter, and county. Take the mean of the AQI value,
  latitude, and longitude.
\item
  Add new columns for ``Month'' and ``Year'' by parsing your ``Date''
  column (hint: \texttt{lubridate} package)
\item
  Hint: the dimensions of this dataset should be 14,752 x 9.
\end{itemize}

\begin{enumerate}
\def\labelenumi{\arabic{enumi}.}
\setcounter{enumi}{9}
\item
  Spread your datasets such that AQI values for ozone and PM2.5 are in
  separate columns. Each location on a specific date should now occupy
  only one row.
\item
  Call up the dimensions of your new tidy dataset.
\item
  Save your processed dataset with the following file name:
  ``EPAair\_O3\_PM25\_NC1819\_Processed.csv''
\end{enumerate}

\begin{Shaded}
\begin{Highlighting}[]
\CommentTok{\#7 see above}

\CommentTok{\#8 see above}

\CommentTok{\#9}

\NormalTok{EPA\_Common\_Sites }\OtherTok{\textless{}{-}}\NormalTok{ Master\_EPA\_Air }\SpecialCharTok{\%\textgreater{}\%} 
\FunctionTok{filter}\NormalTok{(Site\_Name }\SpecialCharTok{\%in\%} \FunctionTok{c}\NormalTok{(}\StringTok{"Linville Falls"}\NormalTok{, }\StringTok{"Durham Armory"}\NormalTok{, }\StringTok{"Leggett"}\NormalTok{, }\StringTok{"Hattie Avenue"}\NormalTok{,  }
    \StringTok{"Clemmons Middle"}\NormalTok{, }\StringTok{"Mendenhall School"}\NormalTok{, }\StringTok{"Frying Pan Mountain"}\NormalTok{, }
    \StringTok{"West Johnston Co."}\NormalTok{, }\StringTok{"Garinger High School"}\NormalTok{, }\StringTok{"Castle Hayne"}\NormalTok{, }
    \StringTok{"Pitt Agri. Center"}\NormalTok{, }\StringTok{"Bryson City"}\NormalTok{, }\StringTok{"Millbrook School"}\NormalTok{)) }\SpecialCharTok{\%\textgreater{}\%} 
      \FunctionTok{mutate}\NormalTok{(}\AttributeTok{month=}\FunctionTok{month}\NormalTok{(Date), }\AttributeTok{year=}\FunctionTok{year}\NormalTok{(Date))}


\NormalTok{EPA\_Summary }\OtherTok{\textless{}{-}}\NormalTok{ EPA\_Common\_Sites }\SpecialCharTok{\%\textgreater{}\%} \FunctionTok{group\_by}\NormalTok{(Date,Site\_Name,AQS\_Parameter\_Desc,County,Site\_Latitude,Site\_Longitude, month, year)}\SpecialCharTok{\%\textgreater{}\%} \FunctionTok{summarise}\NormalTok{(}\AttributeTok{Mean\_AQI =} \FunctionTok{mean}\NormalTok{(Daily\_AQI\_Value))  }\CommentTok{\#IDK if I should use mutate here instead of summarize}
\end{Highlighting}
\end{Shaded}

\begin{verbatim}
## `summarise()` has grouped output by 'Date', 'Site_Name', 'AQS_Parameter_Desc',
## 'County', 'Site_Latitude', 'Site_Longitude', 'month'. You can override using
## the `.groups` argument.
\end{verbatim}

\begin{Shaded}
\begin{Highlighting}[]
  \CommentTok{\#IDK if I should use mutate here instead of summarize}

\CommentTok{\#note to self...at this point the EPA\_Summary dimensions are correct if I use summarise{-} still 14752 9, }
\CommentTok{\#BUT the month and year columns disappear. I added them manually. }


\NormalTok{EPA\_Summary\_AQS\_Adjust }\OtherTok{\textless{}{-}}\NormalTok{EPA\_Summary }\SpecialCharTok{\%\textgreater{}\%} 
  \FunctionTok{pivot\_wider}\NormalTok{(}\AttributeTok{names\_from =}\NormalTok{ AQS\_Parameter\_Desc, }\AttributeTok{values\_from =}\NormalTok{ Mean\_AQI)}
  

\CommentTok{\#This is not quite doing what I want so I think I need to surrender for now. }


\CommentTok{\#11}

\FunctionTok{dim}\NormalTok{(EPA\_Summary\_AQS\_Adjust)}
\end{Highlighting}
\end{Shaded}

\begin{verbatim}
## [1] 8976    9
\end{verbatim}

\begin{Shaded}
\begin{Highlighting}[]
\CommentTok{\#dim are 8976 7....Sounds off.. Not sure where I went wrong}

\CommentTok{\#12}

\FunctionTok{write.csv}\NormalTok{(}
\NormalTok{  EPA\_Summary\_AQS\_Adjust, }
  \AttributeTok{row.names=}\ConstantTok{FALSE}\NormalTok{,}
  \AttributeTok{file=}\FunctionTok{here}\NormalTok{(}\StringTok{\textquotesingle{}./Data/processed/EPAair\_O3\_PM25\_NC1819\_Processed.csv\textquotesingle{}}\NormalTok{)}
\NormalTok{  )}
\end{Highlighting}
\end{Shaded}

\subsection{Generate summary tables}\label{generate-summary-tables}

\begin{enumerate}
\def\labelenumi{\arabic{enumi}.}
\setcounter{enumi}{12}
\item
  Use the split-apply-combine strategy to generate a summary data frame.
  Data should be split into groups by site, month, and year. Then
  compute the mean AQI values for ozone and PM2.5 for each group.
  Finally, add a pipe to remove instances where the mean \textbf{ozone}
  values are not available (use the function \texttt{drop\_na} in your
  pipe). It's ok to have missing mean PM2.5 values in this result.
\item
  Call up the dimensions of the summary dataset.
\end{enumerate}

\begin{Shaded}
\begin{Highlighting}[]
\CommentTok{\#13}

\NormalTok{EPA\_Summaries }\OtherTok{\textless{}{-}} 
\NormalTok{  EPA\_Summary\_AQS\_Adjust }\SpecialCharTok{\%\textgreater{}\%} 
  \FunctionTok{group\_by}\NormalTok{(Site\_Name,month,year,PM2}\FloatTok{.5}\NormalTok{, Ozone) }\SpecialCharTok{\%\textgreater{}\%} 
  \FunctionTok{drop\_na}\NormalTok{(Ozone)}

\CommentTok{\#Question to John: Didn\textquotesingle{}t we already computer the mean AQI in an earlier command?}
\CommentTok{\#I\textquotesingle{}m probably misinterpreting something here.}

\CommentTok{\#14}
\end{Highlighting}
\end{Shaded}

\begin{enumerate}
\def\labelenumi{\arabic{enumi}.}
\setcounter{enumi}{14}
\tightlist
\item
  Why did we use the function \texttt{drop\_na} rather than
  \texttt{na.omit}? Hint: replace \texttt{drop\_na} with
  \texttt{na.omit} in part 13 and observe what happens with the
  dimensions of the summary date frame.
\end{enumerate}

\begin{quote}
Answer: It looks like it deletes objects where the average AQS is
missing/na - \#so the dimensions decrease accordingly.
\end{quote}

\begin{enumerate}
\def\labelenumi{\arabic{enumi}.}
\setcounter{enumi}{15}
\tightlist
\item
  Stage, commit, and push your Assignment to your GitHub account.
  Provide a link to your repository below.
\end{enumerate}

\begin{quote}
Github repository URL: \url{https://github.com/}
\end{quote}

\end{document}
